\subsection{The exact lifting locus on the rank-two elliptic quotient}

\begin{statement}{A.161}{Proved}
{A quadratic cover whose full Selmer lifts have exactly four compact classes}
Let
\[
 \begin{gathered}
 n=1146,\qquad E/\Q:\ y^2=x^3-n^2x,\qquad
 V=E\setminus E[2],\\
 q_0=\frac{M}{2\cdot191^2},\qquad
 C=8\cdot573^2q_0=36M,\qquad M=325596479663.
 \end{gathered}
\tag{A.161.1}
\]
The actual degree-two finite etale cover
\[
 q_{\rm old}:D^\circ\longrightarrow V,
       \qquad x=n(2T-1),\quad y=CW^2
\tag{A.161.2}
\]
is the \(\mu_2\)-torsor \(W^2=y/C\).  Let
\(\mathcal S_0\) be the set of geometric-conjugacy
classes of full Selmer sections \(s\) of \(V/\Q\)
with
\[
                       s^*\kappa_2(y/C)=0
                                  \quad\text{in }H^1(\Q,\mu_2).
\tag{A.161.3}
\]
This is exactly the set of full Selmer sections which
admit a full Selmer lift through (A.161.2).
Every member of \(\mathcal S_0\) is cuspidal.

Compactification and the origin of \(E\) give its
full Tate-module Kummer class \(\eta(s)\).  Its exact
image on this domain is
\[
       \{\eta(s):s\in\mathcal S_0\}
         =\kappa\bigl(E[2](\Q)\bigr)
         =\kappa\{O,(-n,0),(0,0),(n,0)\}
       \subset H^1_{\rm cts}(\Q,TE).
\tag{A.161.4}
\]
Equivalently, the subset of \(\widehat{E(\Q)}\) whose
Kummer images are compact classes of these full Selmer
sections is exactly the embedded group \(E[2](\Q)\).
The point group has rank two by A.143.  Thus the
classification applies to its entire profinite completion,
with no restriction on the number of nonzero primary
coordinates, while retaining the full-section and
local-matching domain in (A.161.3).

The cusp subsets in \(\mathcal S_0\) are explicit.
At any finite two-torsion point use the parameter \(y\)
and its completed tangent class \(\tau_y\); at \(O\)
use \(u=-x/y\) and \(\tau_u\).  Then their lifting
conditions are, respectively,
\[
             [\tau_y]_2=[C]_2=[M]_2,
             \qquad [\tau_u]_2=[-C]_2=[-M]_2.
\tag{A.161.5}
\]
Each of these four subsets is nonempty and is the
image of the corresponding rational cusp packet on
\(D^\circ\).
\end{statement}

\begin{proof}
\emph{The literal quadratic torsor and its inverse.}
The constants satisfy \(a=573q_0^2\) and
\(C^2=8n^3a\).  Therefore
\[
 (CW^2)^2=8n^3T(T-1)(T-\tfrac12)
            =\bigl(n(2T-1)\bigr)^3-n^2\bigl(n(2T-1)\bigr).
\tag{A.161.6}
\]
Conversely a point of \(V\) with a chosen root
\(W^2=y/C\) gives
\[
                            T=\frac{x/n+1}{2}.
\tag{A.161.7}
\]
The equation of \(E\) verifies
\(aW^4=T(T-1)(T-1/2)\).  These two formulas are inverse
on the open curves.  Since \(y\) is nonzero on \(V\),
the derivative \(2W\) is a unit there.  The map is
therefore finite etale of degree two, with deck map
\((T,W)\mapsto(T,-W)\).
Its geometric connectedness follows from the divisor
\[
                \operatorname{div}(y)
                   =(-n,0)+(0,0)+(n,0)-3O:
\tag{A.161.8}
\]
every displayed multiplicity is odd, so \(y/C\) is
not a square even over \(\overline\Q(E)\).
The proper cover has ramification index two exactly
at these four points, and its rational boundary map is
\[
 P_0\mapsto(-n,0),\quad P_{1/2}\mapsto(0,0),\quad
 P_1\mapsto(n,0),\quad P_\infty\mapsto O.
\tag{A.161.9}
\]
The associated geometric character sends each of the
four boundary inertia generators to the nonidentity
element of \(\mu_2\).  Its arithmetic torsor is the
one defined by the actual function \(y/C\), so its
pullback along \(s\) is precisely (A.161.3).

\emph{The full-section lift and every local comparison.}
A zero torsor class in (A.161.3) means that the
section-pulled-back two-element scheme has a
\(G_\Q\)-fixed sheet.  Selecting it places the
section image in the actual open subgroup
\(\pi_1(D^\circ)\subset\pi_1(V)\), after geometric
conjugation.  This constructs a continuous full lift.
Conversely a lift fixes its selected sheet and makes
the class zero.  At a finite ordinary local point,
the fixed sheet is a local rational root of \(y/C\),
and (A.161.7) gives its actual point on \(D^\circ\).
At a cusp the boundary fibre has one point, with
inertia orbit of size two.  The lifted section lies
in its rational decomposition group: its projection
to the full local Galois group makes the residue
field that of the base cusp.  The same argument at
the real place gives the appropriate point or cusp
component.  Thus the constructed lift is Selmer.
These finite-cover maps are compatible with every
further section-neighbourhood, since the selected
subgroup and the restriction of the original full
section define them.

A.160 now puts this lift in one of the four rational
cusp packets of \(D^\circ\).  Under the proper map
(A.161.2), the cusp decomposition group maps into
the cusp decomposition group at its image in
(A.161.9).  Consequently the original \(s\) is
cuspidal.  Compactification kills its boundary
inertia and sends it to the actual point section
at that two-torsion point of \(E\).  Pulling back
the multiplication covers of \(E\) therefore gives
exactly its Kummer class at every level and hence
the full class in (A.161.4).

\emph{The tangent maps and the converse occurrence of all four classes.}
At a finite two-torsion point, \(y\) is a uniformizer:
the derivative \(3x^2-n^2\) is nonzero there.
At its unique preimage in \(D\), \(W\) is a
uniformizer, and the local map is exactly
\[
                                 y=CW^2.
\tag{A.161.10}
\]
At infinity, put \(v=W/T\), a uniformizer on
\(\overline D\), since the orders of \(W,T\) are
\(-3,-4\).  The actual rational functions satisfy
\[
          u=-\frac{x}{y}
            =-\frac{q_0}{2}\,v^2\frac{T}{T-1}.
\tag{A.161.11}
\]
To check this identity, replace \(aW^4\) by
\(T(T-1)(T-1/2)\), and use
\(2na/C=q_0/2\).  The last factor has constant
term one in \(\Q[[v]]\).  Such a formal unit has
a unique \(m\)-th root with constant term one for
every positive integer \(m\), obtained recursively
because \(m\) is invertible in \(\Q\).  It
contributes zero to every tangential Kummer class.
Thus the exact maps on the completed cusp parameters
are
\[
               \tau_y=C\tau_W^2,
               \qquad \tau_u=-\frac{q_0}{2}\tau_v^2.
\tag{A.161.12}
\]
Also \(C/(q_0/2)=(4\cdot573)^2\), so the second
coefficient has squareclass \([-C]\).  These maps
prove (A.161.5).

Conversely every completed class in (A.161.5) has
the indicated square root after division by the
coefficient in (A.161.12).  Explicitly,
\(\widehat{\Q^\times}\) has its sign and completed
prime valuations; reduction modulo two has kernel
exactly the image of multiplication by two.
Dividing the even valuation coordinates by two
preserves their convergence to zero, since
divisibility modulo \(2m\) before division gives
divisibility modulo \(m\) afterwards.  The sign
of an element in that kernel is positive.  This
constructs the completed square root, including
its compatibility over all finite moduli.
It therefore gives the corresponding full cusp
lift through (A.161.12).  Rational nonzero choices
of \(\tau_W\) or \(\tau_v\), for example one,
already supply a full tangential section at each
rational boundary point.  Its pushforward is in
\(\mathcal S_0\).  All four compact classes in
(A.161.4) occur.

Finally the finite Kummer injections
\(E(\Q)/mE(\Q)\hookrightarrow H^1(\Q,E[m])\)
give an injection of \(\widehat{E(\Q)}\) into
the inverse limit of these cohomology groups.
For a completed point, the compatible finite classes
have a compatible system of actual cocycles: the
representatives of each fixed class form a finite
zero-cochain orbit, and representatives at a common
multiple solve every preceding finite set of levels.
Compactness of these finite sets supplies a continuous
\(TE\)-cocycle, defining the stated full Kummer map.
Continuous classes with equal finite components
are equal by the compatible finite-zero-cochain
argument of A.154.  Thus (A.161.4) identifies the
stated subset of the profinite point group exactly.
It does not assign a section to an arbitrary
profinite point: the domain remains the full
Selmer-section set defined by (A.161.3).

The exact torsor, inverse, divisor and tangent
identities are checked in
\path{code/harmonic_positive_rank_lift_certificate.py}
with retained output.
\end{proof}
